% =====================================================================
%  BOZZA - Capitolo 6 (italiano)
%  Conclusioni e lavori futuri
%
%  Risultati in forma qualitativa, con rimandi a \ref{ch:risultati} e
%  placeholder "% TODO numero" da riempire quando il Cap. 5 sara' pronto.
%  La sezione sul gate di forwarding density-aware (lavoro futuro) e'
%  inclusa via \input dallo snippet gia' esistente.
%  Richiede i pacchetti gia' presenti nel preambolo di main.tex.
% =====================================================================
\chapter{Conclusioni e lavori futuri}
\label{ch:conclusioni}

% ---------------------------------------------------------------------
\section{Sintesi dei contributi}
\label{sec:concl-contributi}
Questa tesi ha esteso un simulatore a eventi discreti per il beaconing broadcast
Wi-Fi senza infrastruttura e lo ha usato per studiare insieme due domande: con
quale frequenza una boa di sicurezza marittima deve trasmettere e fino a dove
deve arrivare la consapevolezza che produce. I contributi sono uno scheduler
context-aware (ACAB) che fonde densit\`a, recency del contatto e mobilit\`a; due
meccanismi di propagazione della consapevolezza, append passivo e forwarded
attivo; una suite di metriche che separa l'affidabilit\`a di link dalla
consapevolezza di rete, dalle collisioni e dal carico di canale; e una campagna
sperimentale su canale calibrato sul campo, con popolazione variabile e mobilit\`a
Random Waypoint.

% ---------------------------------------------------------------------
\section{Risultati principali}
\label{sec:concl-risultati}
I risultati completi sono nel Capitolo~\ref{ch:risultati}. Se ne ricavano tre
linee principali.
\begin{enumerate}[leftmargin=1.4em]
  \item \textbf{L'adattamento context-aware costa poco e rende molto.} ACAB tiene
  il tasso di collisione basso dove il beaconing statico cresce molto con la
  densit\`a, e trasmette molti meno beacon sia di SBP sia di ADAB, con una PDR per
  ricevitore pari o migliore a quella delle due linee di base.
  % TODO numero: tasso di collisione ACAB vs SBP alla densita' massima
  % TODO numero: riduzione percentuale dei beacon trasmessi (ACAB vs SBP; ACAB vs ADAB)
  \item \textbf{Affidabilit\`a di link e consapevolezza di rete vanno misurate
  separatamente.} Sotto un canale con path loss realistico la PDR per ricevitore
  \`e limitata dalla propagazione, non dalla contesa, quindi il valore
  dell'adattamento emerge nel tasso di collisione e nel carico di canale pi\`u
  che nella consegna grezza.
  \item \textbf{La consapevolezza passiva batte la propagazione attiva.} La
  modalit\`a append raggiunge una porzione di rete pari o maggiore della
  modalit\`a forwarded, generando molto meno traffico e con una latenza di
  reazione inferiore; il piggyback dei vicini amplifica di parecchio la
  consapevolezza a un hop al solo costo di beacon pi\`u grandi.
  % TODO numero: rapporto di traffico append vs forwarded; fattore di latenza
\end{enumerate}
Alla base delle tre linee c'\`e un principio: in una rete broadcast senza
infrastruttura, l'informazione locale e la parsimonia nell'uso del canale sono
sufficienti, e spesso convengono.

% ---------------------------------------------------------------------
\section{Limitazioni}
\label{sec:concl-limiti}
Il livello fisico \`e semplificato di proposito: probabilit\`a di ricezione
costante a tratti, un solo tasso PHY ($1\,$Mbps), nessun capture effect e nessun
fading esplicito. Il modello di mobilit\`a, pur migliorato a Random Waypoint, non
cattura la deriva correlata dovuta alle correnti n\'e il raggruppamento degli
utenti. La scala provata copre densit\`a costiere realistiche ma non
l'affollamento estremo, e l'analisi del carico di canale considera le
trasmissioni come proxy dell'airtime, senza un budget energetico completo. I pesi
di ACAB $(0{,}4, 0{,}3, 0{,}3)$ sono stati fissati e non ottimizzati, e il
confronto fra append e forwarded \`e stato fatto in un dominio ben connesso.

% ---------------------------------------------------------------------
\section{Lavori futuri}
\label{sec:concl-futuro}
Alcune direzioni seguono in modo naturale. Un modello di perdita dipendente dalla
lunghezza del pacchetto permetterebbe di pesare i beacon crescenti dell'append
contro i frame in pi\`u del forwarded sullo stesso piano, e mostrerebbe se liste
di vicini molto grandi finiscano per danneggiare la consegna. I pesi di ACAB
invitano a una taratura sistematica, eventualmente online. Valutare entrambe le
modalit\`a di consapevolezza in topologie deliberatamente \emph{partizionate}
metterebbe alla prova l'ipotesi che il forwarding ripaghi il suo costo solo
quando il piggyback non pu\`o raggiungere fisicamente i nodi. Tracce GPS reali di
nuotatori e piccoli mezzi affinerebbero il modello di mobilit\`a, e un budget
energetico completo, comprensivo della catena di ricezione e dell'ascolto,
raffinerebbe le stime sui consumi. Infine, il simulatore e i suoi risultati si
trasferiscono oltre il caso marittimo, a qualunque problema di scoperta broadcast
senza infrastruttura, come i beacon di sicurezza veicolari, le reti ad hoc per la
risposta alle emergenze e l'IoT denso.

% Sottosezione "Gate di forwarding density-aware" (lavoro futuro), gia' scritta.
% =====================================================================
%  SNIPPET - Lavori futuri: gate di forwarding density-aware
%
%  Da incollare nel futuro Capitolo 6 (Conclusioni e lavori futuri),
%  come sottosezione. Definisce \label{eq:fwdgate}, riutilizzabile in
%  eventuali rimandi.
%
%  Richiede amsmath nel preambolo.
% =====================================================================

\subsection{Gate di forwarding density-aware}
\label{sec:future-fwdgate}

La modalit\`a forwarded descritta nel Capitolo~\ref{ch:sistema} rilancia ogni
beacon fresco, entro TTL e non proprio, con il solo limite della coda. In reti
dense questo pu\`o produrre molti rilanci ridondanti. Un'estensione possibile \`e
un \emph{gate di forwarding} che decide, una volta sola al momento
dell'accodamento, se rilanciare un beacon ricevuto, con probabilit\`a
\begin{equation}
  P_{\text{fwd}} =
  \begin{cases}
    1 & N \le n_0,\\[2pt]
    \dfrac{n_0}{N}\cdot \max\!\big(0.3,\;1-f_q\big) & N > n_0,
  \end{cases}
  \label{eq:fwdgate}
\end{equation}
dove $N$ \`e il numero di vicini del nodo che rilancia, $n_0$ \`e un numero base
atteso di inoltranti per cascata e $f_q$ \`e il punteggio di backoff (al quadrato)
calcolato per ultimo dallo scheduler. Il fattore $n_0/N$ dirada i rilanci nelle
zone dense, cos\`i che in media solo $\approx n_0$ degli $N$ candidati rilancino.
Il fattore $\max(0.3,\,1-f_q)$ riusa lo stesso segnale di congestione locale che
governa la frequenza dei beacon propri per governare anche i rilanci: una boa che
sta gi\`a rallentando i propri beacon riduce anche i rilanci, fino a un minimo di
$0.3$.

Il gate \`e gi\`a stato progettato ma non \`e attivo nell'implementazione usata
per gli esperimenti di questa tesi. La sua valutazione --- impatto su rilanci,
collisioni, latenza e copertura, e taratura di $n_0$ --- \`e lasciata ai lavori
futuri.

